% ============================================================
% 00_abstract.tex — Abstract (EN + MK)
% ============================================================

\section*{Abstract in English}

This thesis presents the design, implementation, and evaluation of a fully automated team management system for Fantasy Premier League (FPL), built on a three-layer hybrid artificial intelligence architecture. The system integrates machine learning, combinatorial optimisation, and large language models (LLMs) to make autonomous decisions about squad selection, transfers, captain choice, and chip activation throughout an entire season without any manual intervention.

The prediction layer employs four separate LightGBM models — one per player position (goalkeeper, defender, midfielder, forward) — to forecast expected points for each gameweek. Models are trained on six seasons of historical FPL data (2019--20 through 2024--25), comprising 51,044 rows of player-gameweek observations organised into four positional files. Features encompass rolling form windows (last 3 and 5 gameweeks), team-level expected goals and expected goals against (xG/xGA) sourced from Understat, fixture difficulty ratings (FDR), market signals (transfer volume, ownership percentage), and new signing statistics from nine European leagues sourced from FBref. Training follows a walk-forward cross-validation protocol with five temporally ordered folds and exponentially increasing fold weights to emphasise recent seasons.

The optimisation layer formulates squad selection as an Integer Linear Programming (ILP) problem, maximising the sum of predicted points for a 15-player squad subject to FPL constraints: a \pounds 100 million budget, positional structure requirements, a maximum of three players per club, and transfer penalty costs of minus four points per additional transfer beyond the banked allowance.

The intelligence layer consists of seven modules that gather real-time information before each deadline: player injury and availability status from the FPL API, press conference transcripts scraped from Fantasy Football Scout, per-player availability scores (0--100), rotation risk scores (0--100), and LLM-generated strategic recommendations via the Gemini 2.5 Flash API. These signals are converted into multiplicative adjustment factors applied to model predictions before optimisation, directly penalising selections that carry elevated injury or rotation risk.

The system was evaluated across gameweeks 1--28 of the 2025--26 FPL season, achieving \textbf{1,799 total points} (average 64.3 points per gameweek) with zero transfer penalty points incurred. Hyperparameters were tuned through a joint random search over 250 trials followed by 100-trial Bayesian optimisation using Optuna's TPE sampler, with the final configuration (Optuna trial 429) surpassing all prior configurations. Results demonstrate that the intelligence pipeline adds substantial value over pure ML-based prediction, and that LightGBM consistently outperforms XGBoost across all hyperparameter configurations explored — with 14 out of 15 top results belonging to LightGBM.

\bigskip
\noindent\textbf{Keywords:} machine learning, optimisation, fantasy sports, LightGBM, ILP, LLM, LLM agent, Bayesian optimisation, time series, Fantasy Premier League

